\documentclass[11pt]{amsart}
\usepackage[margin=1in]{geometry}
\usepackage{amssymb,mathtools}
\usepackage[hidelinks]{hyperref}

\newtheorem{theorem}{Theorem}
\newtheorem{corollary}[theorem]{Corollary}
\newtheorem{lemma}[theorem]{Lemma}
\newtheorem{proposition}[theorem]{Proposition}

\newcommand{\Z}{\mathbb{Z}}
\newcommand{\Q}{\mathbb{Q}}
\newcommand{\F}{\mathbb{F}}
\newcommand{\cO}{\mathcal{O}}
\newcommand{\cM}{\mathcal{M}}
\newcommand{\sd}{\sqrt{d}}

\title{Landau--Siegel zeros and a Stepanov-type nonvanishing lemma}
\author{GPT-6 Pro and Claude Opus 5.5}
\date{18 October 2026}

\begin{document}

\begin{abstract}
Let $\chi$ be a primitive real Dirichlet character of conductor $q$. We show that
\[
\sum_{p\le X,\ \chi(p)=-1}\frac{\log p}p\le\frac12\log X+\frac12\log q+\log\log X+O(1)
\]
uniformly in $X\ge3$ and $q$; consequently $L(s,\chi)$ has no Landau--Siegel zero. This is a variant
of the OpenAI proof, following Green and Kerr. The key arithmetic observation is that Frobenius at a prime
inert in the quadratic field of $\chi$ is conjugation. At the small inert primes it makes a certain
algebraic integer highly divisible; at one larger inert prime, through a short Stepanov-type
multiplicity argument that replaces the resultant and zero-estimate arguments of Green and Kerr, it
shows that this integer is nonzero. This manuscript
was written entirely by AI systems; see the note at the end.
\end{abstract}

\maketitle

\section{Introduction}

Let $\chi$ be a primitive real Dirichlet character of conductor $q\ge3$, and let $S_\pm(X)$ denote
the sum of $(\log p)/p$ over the primes $p\le X$ with $\chi(p)=\pm1$. Implied constants are
absolute.

\begin{theorem}\label{thm:main}
For every $X\ge3$, \ $S_-(X)\le\frac12\log X+\frac12\log q+\log\log X+O(1)$.
\end{theorem}

\begin{corollary}\label{cor:siegel}
There is an absolute constant $c>0$ such that every zero $\beta\in(0,1)$ of $L(s,\chi)$ satisfies
$(1-\beta)\log q\ge c$.
\end{corollary}

For fixed $q$ one has $S_-(X)=\frac12\log X+O_q(1)$, by Mertens' theorem and
$\sum_{p\le X}\chi(p)(\log p)/p=O_q(1)$ (a consequence of $L(1,\chi)\ne0$). So the coefficient
$\frac12$ of $\log X$ cannot be lowered, and the content of Theorem~\ref{thm:main} is its uniformity
in $q$; a real zero near $1$ would make $S_-(X)$ nearly $\log X$ for $X$ a large
power of $q$ (\S\ref{sec:cor}).

\emph{The idea.} Write $\chi=(d/\cdot)$ with $d$ a fundamental discriminant, $|d|=q$, and for
$\alpha=x+y\sd$ put $\alpha'=x-y\sd$. If $\chi(p)=-1$, conjugation on $\Z[\sd]/p$ is the
Frobenius, so at every point $(z,w)=(\alpha,\alpha')$
\begin{equation}\label{eq:fact}
w\equiv z^p\pmod p,\qquad\text{and}\qquad\partial_z^\kappa(z^p-w)\equiv0\pmod p\quad(\kappa\ge1).
\end{equation}
Let $\delta$ be a nonzero minor of the matrix of $z$-derivatives of monomials $z^iw^j$ at the points
$(\alpha,\alpha')$, $\alpha$ in an $N\times N$ box. At each small inert prime $p$,~\eqref{eq:fact}
makes $\delta$ divisible by about $p^{I/p}$, where $I$ is the total $z$-degree. At one inert prime
$p_0$ above all the $z$-degrees,~\eqref{eq:fact} shows that such a $\delta$ exists: modulo $p_0$ the
points lie on the graph $w=z^{p_0}$ of Frobenius, and a vanishing combination of rows would be a
curve $F=0$ meeting that graph with multiplicity at least $t$ (the number of derivatives) at each
of $p_0$ points, whereas $F(z,z^{p_0})$ has degree below $p_0t$. A box of side $N\approx\sqrt{p_0}$
supplies the $p_0$ points.
Ignoring the costs of the $w$-degrees and of differentiation, $|\delta|$ is then at most about $H^I$
with $H\approx N\sqrt q$, and comparison gives
$\sum_{p\ \mathrm{inert}}(\log p)/p\lesssim\log H\approx\frac12\log p_0+\frac12\log q$. For a split
prime Frobenius is the identity and $\partial_z(z^p-z)\equiv-1$ is a unit, so only $S_-$ is controlled; that is the side a Landau--Siegel
zero inflates.

\emph{Earlier work.} The strategy---evaluating polynomials at Frobenius-twisted points of a box,
and comparing the archimedean size of the resulting algebraic integer with its divisibility by inert
primes---is that of the OpenAI proof~\cite{OpenAI} (over $\Q(\sd,\sqrt2)$), as expounded by
Banks~\cite{Banks}; we follow the variants of Kerr~\cite{Kerr} and Green~\cite{Green}. From Green we
take the field $\Q(\sd)$, the grid of points $(\alpha,\alpha')$, the derivative evaluations, the
height bound, the Leibniz-rule divisibility of the differentiated rows, and (as remarked
in~\cite[\S2]{Green}) the use of a nonsingular minor in place of a Gram determinant; from Kerr, the weighted
monomials and the unitriangular change of basis to $z^v(z^p-w)^h$, which~\cite[Lemma~5.1]{Green}
also uses in truncated, unweighted form. The grid has side about $\sqrt{p_0}$, as in~\cite{Green}
(where $N=2X^{1/2}$).

The main change is the nonvanishing argument (Lemma~\ref{lem:stepanov}): the multiplicity argument at
the heart of Stepanov's method~\cite{Stepanov}, applied at an auxiliary inert prime. It plays the
role of Masser's zero estimate~\cite{Masser} in~\cite{Kerr} and of the resultant argument
of~\cite[Lemma~3.1, Appendix~A]{Green}. Green's lemma is more general: it applies to a much broader
class of point configurations under geometric hypotheses, showing that the nonvanishing is a generic phenomenon and that the
arithmetic lies in the divisibility. Lemma~\ref{lem:stepanov} instead exploits the arithmetic of these
particular points; its merit is brevity, and that it reuses the fact~\eqref{eq:fact} behind the
divisibility. Further, because
Kerr's weighted monomials need no truncation in the change of basis, the weight $\ell$ may grow like
$\log X$; this makes the costs of the $w$-degrees and of differentiation negligible and gives a bound
uniform in $X$ and $q$, a modest sharpening of the bound $(\frac12+O(1/C))\log X+O_C(1)$ for
$X=q^C$ in~\cite[Theorem~1.1]{Green}.

\section{Two lemmas}

Let $\cO=\Z[\sd]=\Z\oplus\Z\sd$.

\begin{lemma}\label{lem:frob}
Let $p$ be an odd prime with $\chi(p)=-1$. Then $\cO/p\cO$ is a field with $p^2$ elements, and
$\alpha^p\equiv\alpha'\pmod{p\cO}$ for every $\alpha\in\cO$.
\end{lemma}

\begin{proof}
As $d$ is not a square modulo $p$, $\cO/p\cO\cong\F_p[T]/(T^2-d)$ is a field; and by Euler's
criterion $(x+y\sd)^p\equiv x+yd^{(p-1)/2}\sd\equiv x-y\sd$.
\end{proof}

\begin{lemma}[after Stepanov~\cite{Stepanov}]\label{lem:stepanov}
Let $k$ be a field of characteristic $p$, let $1\le t\le p$, and let $F\in k[z,w]$ have
$\deg_zF<p$ and $\deg_wF<t$. If $(\partial_z^\kappa F)(a,a^p)=0$ for $0\le\kappa<t$ at $p$ distinct
points $a\in k$, then $F=0$.
\end{lemma}

\begin{proof}
Suppose $F\ne0$ and put $G(z)=F(z,z^p)$. The monomial $z^iw^j$ becomes $z^{i+pj}$, and these
exponents are distinct because $0\le i<p$; so $G\ne0$ and $\deg G\le(p-1)+p(t-1)=pt-1$. As the
derivative of $z^p$ vanishes in characteristic $p$, the chain rule gives
$G^{(\kappa)}(z)=(\partial_z^\kappa F)(z,z^p)$, so $G^{(\kappa)}(a)=0$ for $\kappa<t$ at each of the
$p$ points. Since $\kappa!\ne0$ in $k$ for $\kappa<t\le p$, Taylor expansion at $a$ shows that
$(z-a)^t$ divides $G$. Hence $\deg G\ge pt$, a contradiction.
\end{proof}

Geometrically, in characteristic $p$ the derivative $\partial_z$ differentiates along the graph
$w=z^p$ of Frobenius, and the proof counts the intersections of $F=0$ with that graph.

\section{Proof of Theorem~\ref{thm:main}}

\begin{proposition}\label{prop:ell}
Let $p_0\ge5$ be a prime with $\chi(p_0)=-1$, put $D=p_0-1$, and let $\ell\ge2$ be an integer. Then
\[
S_-(D)\le\Bigl(\frac12+\frac{7}{2\ell}\Bigr)\log D+\frac{\ell+1}{2\ell}\log q+\log\ell+O(1).
\]
\end{proposition}

\begin{proof}[Deduction of Theorem~\ref{thm:main}]
If $X<q$, Mertens' theorem gives $S_-(X)\le\log X+O(1)\le\frac12\log X+\frac12\log q+O(1)$. Let
$X\ge q$, and let $p_0$ be the largest prime $\le X$ with $\chi(p_0)=-1$; if there is none, or
$p_0\le3$, then $S_-(X)=O(1)$. Otherwise $S_-(X)=S_-(p_0-1)+(\log p_0)/p_0$, since by
maximality the primes in $(p_0,X]$ contribute nothing. Apply
Proposition~\ref{prop:ell} with $\ell=\lceil\log X\rceil\ge2$: as $\log D\le\log X$ and
$\log q\le\log X\le\ell$, we have $\frac7{2\ell}\log D\le\frac72$,
$\frac{\ell+1}{2\ell}\log q\le\frac12\log q+\frac12$, and $\log\ell\le\log\log X+\log2$.
\end{proof}

\begin{proof}[Proof of Proposition~\ref{prop:ell}]
\emph{Monomials.} Let $s=\lfloor D/\ell\rfloor$, and let $\cM$ be the set of $(i,j)\in\Z_{\ge0}^2$
with $i+\ell j\le D$; thus $i\le D<p_0$ and $j\le s$. Put $m=|\cM|$, $I=\sum_\cM i$ and
$J=\sum_\cM j$. As in~\cite{Kerr} (where $\ell=8$), trading $z^p$ for $w$ lowers the weight $i+\ell j$ when
$p>\ell$, which keeps the change of basis below inside $\cM$. Reflecting $i\mapsto D-\ell j-i$ for
each fixed $j$ gives $2I+\ell J=Dm$. Writing $i=\ell u+r$ with $0\le r<\ell$, for each $r$ the
pairs $(u,j)$ fill a triangle $u+j\le\lfloor(D-r)/\ell\rfloor$, symmetric in $u$ and $j$; so
$I\ge\ell J$. Therefore
\begin{equation}\label{eq:counts}
J\le I/\ell,\qquad m\le3I/D.
\end{equation}

\emph{A nonzero algebraic integer.} Let $N=\lceil\sqrt{p_0}\,\rceil$, and let $M$ be the matrix over
$\cO$ with rows indexed by $\cM$, ordered by increasing $i$, columns indexed by $(x,y,\kappa)$ with
$0\le x,y<N$ and $0\le\kappa\le s$, and entries
\[
M_{(i,j),(x,y,\kappa)}=\partial_z^\kappa\bigl(z^iw^j\bigr)(\alpha,\alpha')
=(i)_\kappa\,\alpha^{i-\kappa}\alpha'^{\,j}\quad(\kappa\le i),\qquad\alpha=x+y\sd,
\]
and $0$ if $\kappa>i$; here $(i)_\kappa=i(i-1)\cdots(i-\kappa+1)$. Its rows are linearly
independent modulo $p_0$. Indeed, a relation among them over the field $\cO/p_0\cO$ is a polynomial
$F=\sum_\cM c_{ij}z^iw^j$ with $(\partial_z^\kappa F)(\bar\alpha,\bar\alpha^{p_0})=0$ for
$\kappa\le s$ (Lemma~\ref{lem:frob}) at the $N^2\ge p_0$ reductions $\bar\alpha$, which are distinct
because $N<p_0$ and $1,\sd$ are linearly independent over $\F_{p_0}$. As $\deg_zF\le D<p_0$,
$\deg_wF\le s$ and $s+1\le p_0$, Lemma~\ref{lem:stepanov} with $t=s+1$ gives $F=0$. Choose once and
for all an $m\times m$ minor $\delta$ of $M$ that is nonzero modulo $p_0$. Then $\delta\ne0$. Let $\delta'$ be its
image under $\sd\mapsto-\sd$, which is the same minor with $\alpha$ and $\alpha'$ interchanged; then
$n=|\delta\delta'|=|N_{\Q(\sd)/\Q}(\delta)|$ is a positive integer.

\emph{Its size.} Put $H=2N\sqrt q$ (as in~\cite[\S4]{Green}). As $|\alpha|,|\alpha'|\le N(1+\sqrt q)\le H$ and
$(i)_\kappa\le D^s$, every entry in row $(i,j)$ of either matrix has modulus at most
$D^sH^{i+j}$, and Hadamard's inequality gives
\begin{equation}\label{eq:upper}
\log n\le2(I+J)\log H+2ms\log D+m\log m.
\end{equation}

\emph{Its divisibility.} Let $p$ be a prime with $\ell<p\le D$ and $\chi(p)=-1$. For $(i,j)\in\cM$
write $i=ph+v$ with $0\le v<p$, and replace the row of $z^iw^j$ by that of
\[
z^v(z^p-w)^hw^j=z^iw^j+\sum_{\nu=1}^{h}(-1)^\nu\tbinom h\nu z^{i-p\nu}w^{j+\nu}
\]
The monomials in the sum have smaller $z$-degree and smaller weight
$i+\ell j-(p-\ell)\nu$, so they lie in $\cM$. As the rows are ordered by increasing $i$, this adds to each row integer
multiples of rows of smaller $i$, so it does not change $\delta$, which uses every row. By Leibniz's rule,
$\partial_z^\kappa\bigl(z^v(z^p-w)^hw^j\bigr)$ is an integer combination of products each containing
$h$ factors $\partial_z^\lambda(z^p-w)$ with $\lambda\ge0$. Each such factor lies in $p\cO$ at
$(\alpha,\alpha')$: for $\lambda=0$ by Lemma~\ref{lem:frob}; for $1\le\lambda\le p$ because
$\partial_z^\lambda z^p=(p)_\lambda z^{p-\lambda}$ and $p\mid(p)_\lambda$; and for $\lambda>p$ it is
zero. So the
new row of $(i,j)$ lies in $p^h\cO$, and therefore $\delta\in p^{E_p}\cO$; conjugating, $\delta'\in p^{E_p}\cO$, so $p^{2E_p}\mid n$, where
\begin{equation}\label{eq:div}
E_p=\sum_{(i,j)\in\cM}\Bigl\lfloor\frac ip\Bigr\rfloor\ge\frac Ip-m.
\end{equation}

\emph{Comparison.} The coprime $p^{2E_p}$ all divide $n$, so by~\eqref{eq:div} and Chebyshev's bound
$\sum_{p\le D}\log p\ll D$,
\[
2I\sum_{\substack{\ell<p\le D\\ \chi(p)=-1}}\frac{\log p}p\le2\sum_pE_p\log p+O(mD)\le\log n+O(mD).
\]
Insert~\eqref{eq:upper} and divide by $2I$. By~\eqref{eq:counts} and $m\le(D+1)^2$, we have
$(I+J)/I\le1+\frac1\ell$, $ms/I\le\frac3\ell$, and $(m\log m)/I$, $mD/I$ are $O(1)$. The
primes $p\le\ell$ add at most $\log\ell+O(1)$ to $S_-(D)$, by Mertens' theorem; so
\[
S_-(D)\le\log H+\frac{\log H+3\log D}{\ell}+\log\ell+O(1).
\]
Here $\log H$ is the height of the grid, the terms with $1/\ell$ are the costs of the $w$-degrees and
of differentiation, and $\log\ell$ accounts for the small primes set aside. Since $N\le2\sqrt D$, we
have $\log H\le\frac12\log(Dq)+\log4$, and the proposition follows.
\end{proof}

%\clearpage
\section{Proof of Corollary~\ref{cor:siegel}}\label{sec:cor}

This deduction is standard (cf.~\cite[Lemma~1]{Kerr}); we include it for completeness. Let $L(\beta,\chi)=0$ with $0<\beta<1$, let $X\ge3$, and put $\sigma=1+1/\log X\in(1,2]$. Taking
real parts in the formula for $L'/L$ that comes from the Hadamard product~\cite[p.~89]{Davenport},
\[
-\frac{L'}{L}(\sigma,\chi)=\frac12\log\frac q\pi
+\frac12\frac{\Gamma'}{\Gamma}\Bigl(\frac{\sigma+\mathfrak a}2\Bigr)
-\sum_\rho\operatorname{Re}\frac1{\sigma-\rho},\qquad \chi(-1)=(-1)^{\mathfrak a},
\]
where $\rho$ runs over the nontrivial zeros. The gamma term is bounded, and every term of the sum
is positive because $\operatorname{Re}\rho\le1<\sigma$. Keeping only $\rho=\beta$, adding
$-\frac{\zeta'}{\zeta}(\sigma)\le\frac1{\sigma-1}+O(1)$, and using
$\frac1{\sigma-1}-\frac1{\sigma-\beta}\le\frac{1-\beta}{(\sigma-1)^2}=(1-\beta)(\log X)^2$,
\[
\sum_{n\ge1}\frac{\Lambda(n)(1+\chi(n))}{n^\sigma}\le\tfrac12\log q+(1-\beta)(\log X)^2+O(1).
\]
All terms on the left are nonnegative, and a prime $p\le X$ with $\chi(p)=1$ contributes
$2(\log p)p^{-\sigma}\ge2e^{-1}(\log p)/p$. Hence
$S_+(X)\le\frac e2(1-\beta)(\log X)^2+O(\log q)$.
On the other hand, the primes with $\chi(p)=0$ divide $q$, so Mertens' theorem gives
$S_+(X)+S_-(X)\ge\log X-\log q-O(1)$. Combining these with Theorem~\ref{thm:main} and
$\log\log X\le\frac14\log X+O(1)$,
\[
\tfrac14\log X\le C\log q+\tfrac e2(1-\beta)(\log X)^2
\]
for an absolute constant $C$. Taking $X=q^\lambda$ with $\lambda=\max(1,8C)$ gives
$(1-\beta)\log q\ge1/(4e\lambda)$.\qed

\subsection*{Authorship}
This manuscript is AI-generated. Everything in it that is not taken from the cited works---in
particular the use of an auxiliary inert prime and the arrangement of the proof---was found and
written by AI systems (GPT-6 Pro and Claude Opus 5.5), building on~\cite{OpenAI,Banks,Green,Kerr}. The human
who prepared the manuscript contributed no mathematical arguments, and only passed material
between the systems and asked for revisions. The proofs have been checked by several independent
AI reviews, not by a human. The manuscript has not been refereed.

\begin{thebibliography}{9}\footnotesize\raggedright\setlength{\itemsep}{0pt}
\bibitem{Banks} W.\,D. Banks, \emph{On the OpenAI proof that Landau--Siegel zeros do not exist},
preprint, 2026, arXiv:2610.15483.
\bibitem{Davenport} H. Davenport, \emph{Multiplicative Number Theory}, 2nd ed., revised by
H.\,L. Montgomery, Graduate Texts in Mathematics 74, Springer, 1980.
\bibitem{Green} B. Green, \emph{On Landau--Siegel zeros}, preprint, 2026, hexagon:2610.00311v2,
\url{https://hexagonmath.org/abs/2610.00311v2}.
\bibitem{Kerr} B. Kerr, \emph{A simpler proof that there are no Landau--Siegel zeros}, preprint,
2026, arXiv:2610.21099.
\bibitem{Masser} D.\,W. Masser, \emph{On polynomials and exponential polynomials in several complex
variables}, Invent. Math. \textbf{63} (1981), 81--95.
\bibitem{OpenAI} OpenAI, \emph{Uniform exclusion of Landau--Siegel zeros}, OpenAI Math Release
preprint, 1 October 2026.
\bibitem{Stepanov} S.\,A. Stepanov, \emph{On the number of points of a hyperelliptic curve over a
finite prime field}, Math. USSR-Izv. \textbf{3} (1969), 1103--1114.
\end{thebibliography}

\end{document}
